% ==========================================================================
% Engineering Design Document: Laboratory-Scale Parametric psi-Resonator
%
% Gary Thomas Alcock
% Density Field Dynamics Research Programme
% ==========================================================================
\documentclass[12pt,letterpaper]{article}

\usepackage[margin=1in]{geometry}
\usepackage{amsmath,amssymb}
\usepackage{physics}
\usepackage{graphicx}
\usepackage{hyperref}
\usepackage{xcolor}
\usepackage{booktabs}
\usepackage{multirow}
\usepackage{enumitem}
\usepackage{siunitx}
\usepackage{fancyhdr}
\usepackage{tocloft}
\usepackage{titlesec}
\usepackage{longtable}

\pagestyle{fancy}
\fancyhf{}
\fancyhead[L]{\small DFD Parametric $\psi$-Resonator: Engineering Design}
\fancyhead[R]{\small Rev.\ 1.0}
\fancyfoot[C]{\thepage}
\fancyfoot[R]{\small CONFIDENTIAL -- Patent Pending}

\title{\textbf{Engineering Design Document}\\[6pt]
  \Large Laboratory-Scale Parametric $\psi$-Resonator\\[6pt]
  \large System Specification, Subsystem Design, and Test Protocol\\[12pt]
  \normalsize Density Field Dynamics Research Programme\\
  Patent \#2: Parametric $\psi$-Resonators and Metamaterial Amplifiers}

\author{Gary Thomas Alcock\\
  \texttt{gary@gtacompanies.com}}
\date{Revision 1.0 --- \today}

\begin{document}
\maketitle
\thispagestyle{fancy}

\begin{center}
\fbox{\parbox{0.9\textwidth}{\centering\textbf{DOCUMENT STATUS:}
Design Specification for Laboratory Prototype\\
Based on U.S.\ Provisional Patent Application (Oct.\ 11, 2025)\\
and DFD v3.2 Appendix R (EM--$\psi$ Back-Reaction Coupling)}}
\end{center}

\vspace{12pt}
\tableofcontents
\newpage

% ======================================================================
\section{Executive Summary}

This document provides the complete engineering specification for a
laboratory-scale parametric $\psi$-resonator---the first device designed
to deliberately search for electromagnetic back-reaction on the scalar
refractive field $\psi$ predicted by Density Field Dynamics (DFD).

The system consists of:
\begin{enumerate}[noitemsep]
  \item A high-$Q$ superconducting RF (SRF) cavity operated in a
    dual-mode (TE$_{011}$ + TM$_{010}$) superposition;
  \item A $2\omega$ parametric pump system phase-locked to the cavity
    field;
  \item A precision frequency/phase readout system sensitive to
    $\psi$-mode excitation;
  \item A closed-loop control architecture maintaining stable operation
    at a target fraction of the parametric oscillation threshold.
\end{enumerate}

\textbf{Key performance targets:}
\begin{itemize}[noitemsep]
  \item Phase 1 (room-temperature Cu cavity): $|\lambda - 1| \lesssim 10^{-3}$
  \item Phase 2 (SRF Nb cavity): $|\lambda - 1| \lesssim 10^{-8}$
  \item Phase 3 (pulsed high-energy): $|\lambda - 1| \lesssim 10^{-14}$
\end{itemize}

% ======================================================================
\section{System Requirements}

\subsection{Functional Requirements}

\begin{longtable}{p{1.2cm}p{11cm}p{2cm}}
\toprule
\textbf{ID} & \textbf{Requirement} & \textbf{Priority} \\
\midrule
\endfirsthead
FR-01 & The system shall support simultaneous TE$_{011}$ and TM$_{010}$
  cavity modes with controllable relative phase. & Mandatory \\
FR-02 & The system shall generate a pump tone at frequency
  $2\omega \approx 2\Omega_\psi$ with phase locked to the cavity field. & Mandatory \\
FR-03 & The system shall modulate an effective cavity parameter
  (boundary reactance, stored energy) at $2\omega$ with depth
  $m \geq 0.01$. & Mandatory \\
FR-04 & The system shall detect coherent signals at frequency $\Omega_\psi$
  with sensitivity $\Delta\psi \lesssim 10^{-14}$. & Target \\
FR-05 & The system shall provide a null-test mode (pure TM$_{010}$)
  that suppresses $\psi$ coupling by geometry transparency. & Mandatory \\
FR-06 & The system shall allow frequency scanning of the pump through
  $\pm 10\gamma_\psi$ around $2\Omega_\psi$. & Mandatory \\
FR-07 & The system shall allow pump phase reversal ($\phi_p \to \phi_p + \pi$)
  within one cavity time constant. & Mandatory \\
FR-08 & The system shall record all telemetry (cavity field, pump
  parameters, environmental sensors) with timestamps to $< 1~\mu$s. & Mandatory \\
\bottomrule
\end{longtable}

\subsection{Performance Requirements}

\begin{longtable}{p{1.2cm}p{8cm}p{2.5cm}p{2.5cm}}
\toprule
\textbf{ID} & \textbf{Parameter} & \textbf{Phase 1 (Cu)} & \textbf{Phase 2 (Nb)} \\
\midrule
\endfirsthead
PR-01 & Cavity quality factor $Q_0$ & $> 10^4$ & $> 10^{10}$ \\
PR-02 & Stored energy $U_0$ & $0.1$ J & $10$--$100$ J \\
PR-03 & Modulation depth $m$ & $0.01$--$0.3$ & $0.01$--$0.1$ \\
PR-04 & Pump phase noise (at 1 Hz offset) & $< 10$ mrad & $< 0.1$ mrad \\
PR-05 & $\psi$-mode frequency resolution & $< 1$ Hz & $< 0.01$ Hz \\
PR-06 & Measurement integration time & $10^3$ s & $10^4$--$10^6$ s \\
PR-07 & Cavity frequency stability & $\delta f/f < 10^{-8}$ & $\delta f/f < 10^{-12}$ \\
PR-08 & Environmental isolation (vibration) & $< 10^{-4}$ g & $< 10^{-7}$ g \\
\bottomrule
\end{longtable}

% ======================================================================
\section{System Architecture}

\subsection{Block Diagram}

The system comprises five major subsystems:

\begin{verbatim}
+------------------+     +------------------+     +------------------+
|  Reference       |---->|  PLL / Phase     |---->|  Pump Generator  |
|  Oscillator      |     |  Detector        |     |  (2w with AGC)   |
+------------------+     +------------------+     +--------+---------+
                                |                          |
                                |                          v
+------------------+     +------+--+--------+     +--------+---------+
|  Data            |<----|  Sensing         |<----|  Resonator /     |
|  Acquisition     |     |  (pickoffs, env, |     |  Metamaterial    |
|  & Analysis      |     |  EM monitors)    |     |  (parametric     |
+------------------+     +------------------+     |   psi gain)      |
                                                  +------------------+
                                                         |
                                                         v
                                                  +------------------+
                                                  |  psi-mode        |
                                                  |  output /        |
                                                  |  detection       |
                                                  +------------------+
\end{verbatim}

\subsection{Subsystem overview}

\begin{enumerate}
  \item \textbf{Cavity Subsystem (CS):} SRF cavity, cryostat, vacuum,
    mode couplers
  \item \textbf{RF Drive Subsystem (RDS):} Fundamental drive, $2\omega$
    pump, amplitude/phase control
  \item \textbf{Control Subsystem (CTS):} PLL, AGC, mode selection,
    interlock
  \item \textbf{Detection Subsystem (DS):} $\psi$-mode readout,
    frequency discriminator, lock-in amplifier
  \item \textbf{Data \& Environmental Subsystem (DES):} DAQ, timing,
    temperature/vibration/magnetic sensors
\end{enumerate}

% ======================================================================
\section{Cavity Subsystem Design}

\subsection{Phase 1: Room-Temperature Copper Cavity}

\subsubsection{Geometry}

\begin{tabular}{ll}
\toprule
Parameter & Value \\
\midrule
Shape & Cylindrical pillbox with coupling irises \\
Material & OFHC copper, electropolished \\
Inner radius $R$ & 150 mm \\
Length $L$ & 300 mm \\
Wall thickness & 5 mm \\
TM$_{010}$ frequency & $f_{010} = 2.405c/(2\pi R) = 766$ MHz \\
TE$_{011}$ frequency & Matched to $f_{010}$ via slight ellipticity \\
 & ($R_a/R_b = 1.015 \pm 0.001$) \\
Volume & $V = \pi R^2 L = 21.2$ L \\
Surface area & $A_s = 2\pi R(R + L) = 0.424$ m$^2$ \\
\bottomrule
\end{tabular}

\subsubsection{Mode matching}

The TE$_{011}$ and TM$_{010}$ modes must be frequency-degenerate.
For a cylinder of radius $R$ and length $L$:
\begin{align}
  f_{\text{TM}_{010}} &= \frac{2.405\,c}{2\pi R} = 766~\text{MHz}, \\
  f_{\text{TE}_{011}} &= \frac{c}{2\pi}\sqrt{\left(\frac{3.832}{R}\right)^2
    + \left(\frac{\pi}{L}\right)^2}.
\end{align}

Setting $f_{\text{TE}_{011}} = f_{\text{TM}_{010}}$ determines $L/R$:
\begin{equation}
  \frac{L}{R} = \frac{\pi}{\sqrt{(2.405)^2 - (3.832)^2}} \cdot R.
\end{equation}

Since $2.405 < 3.832$, exact degeneracy in a circular cylinder is not
possible for TE$_{011}$/TM$_{010}$.  Instead, we use:
\begin{itemize}[noitemsep]
  \item \textbf{Option A:} Slightly elliptical cross-section to split
    the TE$_{011}$ into two polarizations, one of which can be tuned
    to match TM$_{010}$ via mechanical deformation (tuner);
  \item \textbf{Option B:} Use TM$_{010}$ + TE$_{111}$ (next higher
    TE mode), with $L$ chosen for degeneracy:
    \begin{equation}
      L = \frac{\pi R}{\sqrt{(1.841)^2 - (2.405)^2 + (\pi R/L_0)^2}}.
    \end{equation}
  \item \textbf{Option C (preferred):} Operate at a higher-order mode
    pair that is naturally degenerate, or use two coupled cavities
    with a slot coupling.
\end{itemize}

For the Phase~1 prototype, we adopt \textbf{Option A} with a
mechanical tuner providing $\pm 5$~MHz tuning range.

\subsubsection{Quality factor}

At 766~MHz, room-temperature copper:
\begin{equation}
  R_s = \sqrt{\frac{\pi f \mu_0}{\sigma_{\text{Cu}}}}
  \approx 7.8~\text{m}\Omega,
\end{equation}
giving:
\begin{equation}
  Q_0 = \frac{G}{R_s} \approx \frac{260~\Omega}{7.8\times 10^{-3}~\Omega}
  \approx 33{,}000,
\end{equation}
where $G \approx 260~\Omega$ is the geometry factor for TM$_{010}$.

With coupling ($\beta = 1$): $Q_L \approx Q_0/2 \approx 16{,}500$.

\subsubsection{Stored energy}

At $E_{\text{peak}} = 2$~MV/m (conservative for Cu):
\begin{equation}
  U_0 = \frac{\varepsilon_0}{2}\int |E|^2\,dV
  \approx \frac{\varepsilon_0}{2}\,\frac{E_{\text{peak}}^2}{2}\,V
  \approx 0.094~\text{J} \approx 0.1~\text{J}.
\end{equation}

\subsection{Phase 2: Superconducting Niobium Cavity}

\begin{tabular}{ll}
\toprule
Parameter & Value \\
\midrule
Material & Bulk niobium, RRR $> 300$, BCP or EP finish \\
Inner radius $R$ & 150 mm \\
Length $L$ & 300 mm (single cell) or 600 mm (2-cell) \\
Wall thickness & 3 mm \\
Operating temperature & 2.0 K (superfluid He) \\
$Q_0$ (2 K, low field) & $> 2 \times 10^{10}$ \\
$Q_0$ (2 K, 40 MV/m) & $> 10^{10}$ (with N-doping) \\
Max.\ $E_{\text{peak}}$ & 50 MV/m \\
Max.\ stored energy & $\sim 200$ J (at 50 MV/m in 2-cell) \\
Operating point & 100 J (conservative, 40 MV/m) \\
Cavity bandwidth & $f_0/Q_L \approx 0.05$ Hz (at $Q_L = 10^{10}$) \\
\bottomrule
\end{tabular}

\subsubsection{Cryostat requirements}

\begin{itemize}[noitemsep]
  \item LHe capacity: $> 500$ L (1 week hold time);
  \item Static heat load: $< 5$ W at 2 K;
  \item Dynamic heat load (RF): $P_{\text{diss}} = U_0\omega/Q_0
    \approx 100 \times 4.8\times 10^9 / 2\times 10^{10} = 24$ W;
  \item Vibration isolation: active/passive to $< 10^{-7}$~g
    at the cavity;
  \item Magnetic shielding: $\mu$-metal + Nb shield,
    residual $B < 5$~nT.
\end{itemize}

% ======================================================================
\section{RF Drive Subsystem}

\subsection{Fundamental drive ($\omega$)}

\begin{tabular}{ll}
\toprule
Parameter & Specification \\
\midrule
Frequency & 766 MHz $\pm$ 50 kHz (tunable) \\
Power (Phase 1, Cu) & 50 W CW \\
Power (Phase 2, Nb) & 200 W CW (klystron or solid-state) \\
Phase noise & $< -130$ dBc/Hz at 1 Hz offset \\
Amplitude stability & $\delta P/P < 10^{-4}$ \\
Frequency reference & 10 MHz from GPS-disciplined Rb standard \\
\bottomrule
\end{tabular}

\subsection{$2\omega$ pump drive}

\begin{tabular}{ll}
\toprule
Parameter & Specification \\
\midrule
Frequency & 1532 MHz $\pm$ 100 kHz \\
Power (Phase 1) & 10 W CW \\
Power (Phase 2) & 100 W CW \\
Phase control & Digital, 16-bit, resolution $< 0.1$ mrad \\
Phase noise & $< -140$ dBc/Hz at 1 Hz offset \\
Modulation depth $m$ & Adjustable 0.001--0.3 \\
Coupling & Dedicated port, variable $Q_{\text{ext}}$ ($10^4$--$10^8$) \\
\bottomrule
\end{tabular}

\subsubsection{Pump coupling scheme}

The $2\omega$ pump tone is injected through a dedicated coupling port
positioned at an EM field maximum for the $2\omega$ cavity mode.
The coupling is via a loop or probe coupler with:
\begin{itemize}[noitemsep]
  \item Adjustable insertion depth (motorized, 0.1~mm resolution);
  \item Directional coupler for reflected power monitoring;
  \item Band-pass filter centered at $2\omega$ to prevent pump leakage
    into the fundamental drive chain.
\end{itemize}

\subsection{Phase-locked loop (PLL)}

\subsubsection{Architecture}

\begin{verbatim}
Cavity pickup     +--------+   +--------+   +--------+
  at omega  ----->| x2     |-->| Phase  |-->| Loop   |
                  | Freq.  |   | Det.   |   | Filter |
                  | Doubler|   | (DDS)  |   | (PID)  |
                  +--------+   +---+----+   +---+----+
                                   |            |
                 2w pump     +-----+-----+  +---+----+
                 <-----------| VCO /     |<-| DAC    |
                             | DDS       |  |        |
                             +-----------+  +--------+
\end{verbatim}

\subsubsection{Specifications}

\begin{tabular}{ll}
\toprule
Parameter & Value \\
\midrule
Phase detector type & Digital (DDS-based, Analog Devices AD9959) \\
Loop bandwidth & 10 kHz (adjustable 100 Hz -- 100 kHz) \\
Phase noise floor & $-150$ dBc/Hz at $> 10$ kHz offset \\
Lock range & $\pm 1$ MHz \\
Capture range & $\pm 100$ kHz \\
Phase resolution & 0.05 mrad (20-bit DAC) \\
Lock indicator & SNR $> 20$ dB in 1 Hz BW \\
\bottomrule
\end{tabular}

\subsection{Automatic gain control (AGC)}

\begin{tabular}{ll}
\toprule
Parameter & Value \\
\midrule
Sensor & Transmitted power from cavity pickup \\
Setpoint & $U_0 = $ programmable (0.01--200 J) \\
Bandwidth & 1 kHz \\
Dynamic range & 60 dB \\
Stability & $\delta U_0/U_0 < 10^{-4}$ over 1000 s \\
\bottomrule
\end{tabular}

% ======================================================================
\section{Detection Subsystem}

\subsection{$\psi$-mode frequency discriminator}

The primary detection channel monitors the cavity resonance frequency
for perturbations at $\Omega_\psi$.  A Pound--Drever--Hall (PDH) lock
provides:
\begin{equation}
  \frac{\delta f}{f} = -\frac{\delta\psi}{2}
  \approx -\frac{u(z_0)\,|q|}{2}.
\end{equation}

\subsubsection{PDH specifications}

\begin{tabular}{ll}
\toprule
Parameter & Value \\
\midrule
Modulation frequency & 10 MHz (sideband) \\
Frequency sensitivity & $10^{-16}/\sqrt{\text{Hz}}$ (Phase 2) \\
Detection bandwidth & DC -- 100 kHz \\
Reference & Ultra-stable laser comb (Phase 2) \\
ADC resolution & 24-bit, 1 MSPS \\
\bottomrule
\end{tabular}

\subsection{Lock-in amplifier chain}

A dual-phase lock-in amplifier referenced to $\Omega_\psi$ extracts
the coherent $\psi$-mode signal from the PDH error signal:
\begin{itemize}[noitemsep]
  \item Reference: derived from the $2\omega$ pump divided by 2;
  \item Time constant: adjustable $0.1$--$10^4$~s;
  \item Noise floor: $1~\text{nV}/\sqrt{\text{Hz}}$;
  \item Outputs: in-phase (X) and quadrature (Y) at $\Omega_\psi$.
\end{itemize}

\subsection{Secondary detection channels}

\begin{enumerate}[noitemsep]
  \item \textbf{Power spectral density:} FFT of cavity transmitted
    power, monitoring for spectral features at $\Omega_\psi$ and
    harmonics;
  \item \textbf{Amplitude modulation:} Envelope detector on cavity
    field, looking for AM at $\Omega_\psi$;
  \item \textbf{Inter-mode beating:} If TE$_{011}$ and TM$_{010}$
    are not exactly degenerate, monitor the beat note for
    $\psi$-induced shifts;
  \item \textbf{External accelerometer:} Mechanical response of the
    cavity/cryostat to any $\psi$ gradient ($\mathbf{a} = (c^2/2)\nabla\psi$).
\end{enumerate}

% ======================================================================
\section{Cavity Geometry and Overlap Integrals}

\subsection{Geometry transparency and its breaking}

For a pure eigenmode in a symmetric cavity, the driven overlap integral
vanishes:
\begin{equation}
  G = \int u(\mathbf{r})\,\bar{\Xi}_{2\omega}(\mathbf{r})\,d^3r
  \approx 0
  \qquad\text{(pure mode, symmetric cavity)}.
\end{equation}

Three methods restore $G \neq 0$:

\begin{enumerate}
  \item \textbf{TE+TM superposition:} Co-phased modes with matched
    radii give:
    \begin{equation}
      G = u(z_0)\,e^{-2\psi_0}\,\eta_\times\,U_0\cos\phi,
    \end{equation}
    where $\eta_\times = \mathcal{O}(0.1$--$1)$ is the cross-mode
    overlap factor and $\phi$ is the relative phase.

  \item \textbf{Asymmetric geometry:} Irises, cutoff asymmetries, or
    conical sections break equipartition.

  \item \textbf{Mode beating:} Two nearby modes at $\omega_1$, $\omega_2$
    produce $2\omega = \omega_1 + \omega_2$ components automatically.
\end{enumerate}

\subsection{Parametric overlap $H$}

For $N$ cavities of total aperture $A_{\text{cav,tot}}$ arranged along
a $\psi$-mode tube of cross-section $A_\psi$:
\begin{equation}
  H \approx \frac{2}{L}\,\kappa_{\text{eff}}\,
    \frac{A_{\text{cav,tot}}}{A_\psi}.
\end{equation}

\subsection{Design parameters for sensitivity calculation}

\begin{center}
\begin{tabular}{lcc}
\toprule
Parameter & Phase 1 (Cu) & Phase 2 (Nb) \\
\midrule
$U_0$ (J) & 0.1 & 100 \\
$m$ & 0.1 & 0.1 \\
$A_{\text{cav,tot}}$ (m$^2$) & $3\times 10^{-3}$ & $3\times 10^{-2}$ \\
$A_\psi$ (m$^2$) & 0.8 & 0.27 \\
$\gamma_\psi/\Omega_\psi$ & $10^{-3}$ & $10^{-3}$ \\
$\kappa_{\text{eff}}$ & 1 & 1 \\
$c_s$ (m/s) & $c$ & $c$ \\
\midrule
$|\lambda - 1|_{\text{min}}$ & $\sim 3\times 10^{-1}$ & $\sim 3\times 10^{-7}$ \\
\bottomrule
\end{tabular}
\end{center}

The Phase~2 Nb cavity substantially exceeds the accidental bound of
$3\times 10^{-5}$ from passive cavity stability, enabling a genuine
discovery search.

% ======================================================================
\section{Metamaterial Amplifier Panel Design}

\subsection{Panel specifications}

\begin{tabular}{ll}
\toprule
Parameter & Value \\
\midrule
Array size & $16 \times 16$ elements \\
Unit cell spacing & $d = \lambda_{\text{EM}}/5 \approx 25$ mm (at 2.4 GHz) \\
Panel dimensions & $400 \times 400$ mm \\
Unit cell type & Split-ring resonator + varactor \\
Varactor & GaAs hyperabrupt, $C = 0.5$--5 pF \\
SRR outer radius & 5 mm \\
SRR gap & 0.5 mm \\
Trace width & 1 mm \\
Substrate & Rogers RT/duroid 5880 ($\varepsilon_r = 2.2$, $\tan\delta = 0.0009$) \\
Substrate thickness & 1.575 mm \\
\bottomrule
\end{tabular}

\subsection{Pump distribution network}

\begin{itemize}[noitemsep]
  \item Corporate feed network distributing $2\omega$ pump to all 256 cells;
  \item Per-cell digital phase shifter (5-bit, $11.25^\circ$ resolution);
  \item Pump power per cell: 10--100 mW;
  \item Total pump power: 2.5--25 W;
  \item Phase gradient for $20^\circ$ steering:
    $\Delta\phi = 2\pi d\sin(20^\circ)/\lambda_{2\omega} \approx 32^\circ$
    per cell.
\end{itemize}

\subsection{Expected $\psi$ gain}

Per-cell stored energy: $U_{\text{cell}} \sim 10^{-9}$ J at
10 mW drive with $Q \sim 200$:
\begin{equation}
  U_{\text{cell}} = \frac{P_{\text{cell}}\,Q}{2\pi f}
  \approx \frac{0.01 \times 200}{2\pi \times 2.4\times 10^9}
  \approx 1.3\times 10^{-13}~\text{J}.
\end{equation}

For the full $16\times 16$ array with coherent addition along
16 rows:
\begin{equation}
  G_{\psi}^{\text{array}} \approx 16 \times g_{\text{cell}}
  \approx 9~\text{dB},
\end{equation}
matching the patent Example~B specification.

% ======================================================================
\section{Expected $\psi$ Gain and Signal Levels}

\subsection{Driven resonance signal ($2\omega = \Omega_\psi$)}

Using Eq.\ (R13) from DFD v3.2 Appendix~R, for TE+TM superposition
with $\eta_\times \sim 0.3$:
\begin{equation}
  \Delta\psi = u(z_0)\,|q|_{\text{res}}
  \approx \frac{|\lambda-1|\,\eta_\times\,U_0\,c_s}{\pi A_\psi\,\gamma_\psi}.
\end{equation}

For Phase~2 ($U_0 = 100$ J, $A_\psi = 0.27$ m$^2$,
$\gamma_\psi = 0.03$ s$^{-1}$):
\begin{equation}
  \Delta\psi \approx 1.2\times 10^{-3}\,|\lambda - 1|.
\end{equation}

This crosses the cavity--atom sensitivity threshold
($\Delta\psi \sim 10^{-15}$) for $|\lambda - 1| \sim 10^{-12}$.

\subsection{Parametric amplification signal ($2\omega \simeq 2\Omega_\psi$)}

The parametric growth rate:
\begin{equation}
  \Gamma = \frac{1}{2}\,h\,\Omega_\psi - \gamma_\psi,
\end{equation}
with $h = (\lambda - 1)\,U_0\,H\,m/(\Mpsi\Opsi^2)$.

Above threshold, the $\psi$ amplitude grows exponentially until
saturation.  The time to reach a detectable level from thermal noise:
\begin{equation}
  t_{\text{detect}} = \frac{1}{\Gamma}\ln\!\left(
    \frac{\Delta\psi_{\text{detect}}}{\Delta\psi_{\text{thermal}}}
  \right).
\end{equation}

For Phase~2, if $|\lambda - 1| = 10^{-6}$ (just above threshold):
$\Gamma \sim 10^{-3}$ s$^{-1}$, $t_{\text{detect}} \sim 10^4$ s
$\approx 3$ hours.

% ======================================================================
\section{Measurement Protocol}

\subsection{Phase 1 protocol (Room-temperature Cu cavity)}

\begin{enumerate}
  \item \textbf{Cavity characterization (1 week):}
    \begin{itemize}[noitemsep]
      \item Map all modes from 600--2000 MHz;
      \item Measure $Q_0$, coupling $\beta$, and field flatness;
      \item Verify TE+TM frequency matching (tune to $< 1$ bandwidth).
    \end{itemize}

  \item \textbf{PLL commissioning (1 week):}
    \begin{itemize}[noitemsep]
      \item Lock $2\omega$ pump to cavity field;
      \item Verify phase noise specification;
      \item Characterize AGC stability.
    \end{itemize}

  \item \textbf{Baseline measurement (1 week):}
    \begin{itemize}[noitemsep]
      \item Pure TM$_{010}$ mode, pump ON: record $10^4$ s;
      \item Pure TM$_{010}$ mode, pump OFF: record $10^4$ s;
      \item Compute difference spectrum at $\Omega_\psi$;
      \item Establish noise floor.
    \end{itemize}

  \item \textbf{Search measurement (2 weeks):}
    \begin{itemize}[noitemsep]
      \item TE+TM superposition, pump ON: record $10^4$ s per frequency step;
      \item Frequency scan: 100 steps across $\pm 10\gamma_\psi$;
      \item Phase reversal at each candidate signal;
      \item Power scaling: vary $U_0$ over 1 decade.
    \end{itemize}

  \item \textbf{Systematics (1 week):}
    \begin{itemize}[noitemsep]
      \item Repeat with cavity rotated $90^\circ$;
      \item Repeat with swapped TE/TM coupling ports;
      \item Vary cryostat thermal conditions;
      \item Document all environmental correlations.
    \end{itemize}
\end{enumerate}

\subsection{Phase 2 protocol (SRF Nb cavity)}

Same protocol structure as Phase~1 with:
\begin{itemize}[noitemsep]
  \item $10\times$ longer integration times (weeks per measurement);
  \item Additional monitoring of microphonics, Lorentz detuning;
  \item Cross-correlation with seismic, thermal, and magnetic channels;
  \item Blind analysis: $\psi$-mode channel processed independently of
    operator knowledge of pump ON/OFF state.
\end{itemize}

\subsection{Decision criteria}

A \textbf{discovery} ($\lambda \neq 1$) requires:
\begin{enumerate}[noitemsep]
  \item Signal-to-noise ratio $> 5\sigma$ at $\Omega_\psi$ in the
    TE+TM configuration with pump ON;
  \item Signal absent (below $2\sigma$) in TM$_{010}$-only null channel;
  \item Signal reverses sign with pump phase reversal;
  \item Signal scales as predicted with $U_0$ and $m$;
  \item Signal survives all systematic cross-checks.
\end{enumerate}

A \textbf{null result} establishes an upper bound:
\begin{equation}
  |\lambda - 1| < \frac{\text{noise floor at } \Omega_\psi}
    {\text{predicted signal per unit } |\lambda - 1|}.
\end{equation}

% ======================================================================
\section{Risk Assessment and Mitigation}

\begin{longtable}{p{3.5cm}p{4cm}p{2cm}p{4.5cm}}
\toprule
\textbf{Risk} & \textbf{Impact} & \textbf{Likelihood} & \textbf{Mitigation} \\
\midrule
\endfirsthead
TE+TM frequency mismatch & Reduced overlap $G$ & Medium &
  Mechanical tuner with 5 MHz range; elliptical cavity backup \\
$2\omega$ pump leakage & False positive at $\Omega_\psi$ & High &
  Band-stop filters; null channel; phase reversal test \\
Microphonics & Noise at $\Omega_\psi$ & High (SRF) &
  Active vibration isolation; correlation analysis \\
Lorentz detuning & Unstable stored energy & Medium (SRF) &
  Fast AGC ($> 1$ kHz); piezo tuner \\
Thermal EMF & Slow drifts & Low &
  Superconducting joints; temperature stabilization \\
$Q_0$ degradation & Reduced sensitivity & Medium &
  Surface processing; clean assembly; N-doping \\
$\psi$-mode damping unknown & $\gamma_\psi$ may be large & High &
  Scan over wide frequency range; multiple cavity sizes \\
$\lambda = 1$ exactly & No signal & Possible &
  Report null bound; motivate matter-based searches \\
\bottomrule
\end{longtable}

% ======================================================================
\section{Budget Estimates}

\subsection{Phase 1 (Room-temperature Cu cavity)}

\begin{tabular}{lr}
\toprule
Item & Est.\ Cost (USD) \\
\midrule
Copper cavity (machined + polished) & \$15,000 \\
RF amplifiers ($\omega$ and $2\omega$) & \$25,000 \\
PLL / DDS electronics & \$10,000 \\
Lock-in amplifier (Stanford SR860) & \$12,000 \\
Frequency reference (Rb + GPS) & \$5,000 \\
DAQ system (NI PXIe) & \$20,000 \\
Cables, connectors, filters & \$5,000 \\
\midrule
\textbf{Phase 1 Total} & \textbf{\$92,000} \\
\bottomrule
\end{tabular}

\subsection{Phase 2 (SRF Nb cavity)}

\begin{tabular}{lr}
\toprule
Item & Est.\ Cost (USD) \\
\midrule
Nb cavity (fabrication, processing) & \$150,000 \\
Cryostat (vertical test stand) & \$200,000 \\
LHe (initial + 1 year operation) & \$50,000 \\
High-power RF (200 W, 766 MHz) & \$80,000 \\
$2\omega$ pump system (100 W, 1.53 GHz) & \$40,000 \\
Digital LLRF system & \$60,000 \\
Vibration isolation platform & \$30,000 \\
Magnetic shielding & \$20,000 \\
Ultra-stable reference (optical comb) & \$100,000 \\
\midrule
\textbf{Phase 2 Total} & \textbf{\$730,000} \\
\bottomrule
\end{tabular}

\subsection{Phase 3 (Pulsed high-energy)}

Phase~3 requires access to an existing particle accelerator facility
(e.g., DESY, Fermilab, Jefferson Lab) with:
\begin{itemize}[noitemsep]
  \item Multi-cell SRF cavity with $U_0 > 10^5$ J (pulsed);
  \item Custom $2\omega$ insertion and detection hardware;
  \item Estimated incremental cost: \$200,000--\$500,000
    (leveraging existing infrastructure).
\end{itemize}

% ======================================================================
\section{Schedule}

\begin{center}
\begin{tabular}{lll}
\toprule
Phase & Duration & Milestone \\
\midrule
Phase 1 design & Months 1--3 & Cavity + electronics ready \\
Phase 1 commission & Months 4--5 & PLL locked, baseline taken \\
Phase 1 search & Months 6--8 & $|\lambda-1| < 10^{-3}$ or discovery \\
Phase 2 cavity fab & Months 6--12 & Nb cavity delivered \\
Phase 2 cryotest & Months 13--15 & $Q > 10^{10}$ verified \\
Phase 2 search & Months 16--24 & $|\lambda-1| < 10^{-8}$ or discovery \\
Phase 3 proposal & Month 12 & Accelerator facility MOU \\
Phase 3 execution & Months 24--36 & $|\lambda-1| < 10^{-14}$ or discovery \\
Metamaterial panel & Months 6--18 & 16$\times$16 panel tested \\
\bottomrule
\end{tabular}
\end{center}

% ======================================================================
\section{Compliance and Safety}

\begin{itemize}[noitemsep]
  \item RF safety: compliance with IEEE C95.1 / ICNIRP guidelines;
    all RF contained within shielded enclosure;
  \item Cryogenic safety: oxygen-deficiency hazards from LHe boil-off;
    ODH monitoring per Fermilab ES\&H guidelines;
  \item Electrical safety: all high-voltage RF components properly
    enclosed and interlocked;
  \item Magnetic fields: Nb cavity stray fields negligible;
    external solenoids (if used) within occupational limits;
  \item Radiation: no ionizing radiation sources; no beam operations.
\end{itemize}

% ======================================================================
\section{Conclusion}

This engineering design document establishes that a laboratory-scale
parametric $\psi$-resonator is fully achievable with current technology.
The Phase~1 copper-cavity prototype requires approximately \$92,000 and
8 months, delivering a first constraint on $|\lambda - 1|$ at the
$10^{-3}$ level.  The Phase~2 SRF system reaches $10^{-8}$, penetrating
well below the accidental bound.  Phase~3, leveraging existing
accelerator infrastructure, achieves the ultimate $10^{-14}$ target.

The experiment has a clear binary outcome: either $\lambda \neq 1$ is
discovered---opening the door to controllable laboratory $\psi$-field
technology---or the bound is pushed to extraordinary precision,
definitively establishing that EM back-reaction on $\psi$ is negligible
at accessible energies.  Either result is of fundamental significance.

\vfill
\begin{center}
\rule{0.5\textwidth}{0.4pt}\\[6pt]
\textit{End of Engineering Design Document}\\
\textit{Revision 1.0}
\end{center}

\end{document}
